Este capítulo analisa os trabalhos relacionados com três problemas do
apoio à decisão no basquetebol: prever resultados, avaliar e selecionar
alinhamentos e planear sequências de decisões. A previsão enquadra-se na
componente de AAUTIA; a otimização e o planeamento, na de PLANTIA. A
avaliação dos alinhamentos estabelece a ligação entre ambas, ao fornecer
estimativas que podem orientar a escolha de ações. A revisão termina com
uma comparação dos trabalhos e das oportunidades de integração identificadas.

\section{Previsão de resultados e explicabilidade}
\label{aautia:estado-da-arte-de-aautia}

A previsão pré-jogo procura estimar o resultado de uma partida com base
na informação disponível antes do seu início. Neste contexto, a escolha
dos algoritmos está ligada à construção das características e ao modo
como se avalia a capacidade de generalização para jogos futuros.
Lampis et al.~\cite{lampis2023} e Horvat et al.~\cite{horvat2023}
exploram estas dimensões em competições europeias e na NBA,
respetivamente, oferecendo referências complementares para a preparação
dos dados e a comparação de modelos.

Lampis et al.~\cite{lampis2023} analisaram 5\,214 jogos de quatro
competições europeias, realizados entre 2013 e 2018. Compararam Regressão
Logística, Random Forest, XGBoost e combinações de modelos, utilizando
estatísticas históricas e sistemas de classificação como Elo, PageRank
e \emph{pi-rating}. Os indicadores adicionais melhoraram a exatidão em
cerca de três a cinco pontos percentuais face às referências mais
simples. Contudo, nenhum algoritmo foi universalmente superior: o
desempenho variou com a competição e as características dos dados.

Horvat et al.~\cite{horvat2023} dão particular atenção à definição
matemática das estatísticas e dos indicadores derivados utilizados na
previsão de jogos da NBA. Em conjunto, estes trabalhos sustentam uma
comparação que avalie tanto o algoritmo como a informação que o alimenta.
Indicadores de eficiência ofensiva e defensiva, forma recente,
lançamentos, ressaltos e perdas de bola constituem possíveis entradas,
desde que calculados apenas com dados disponíveis no momento da previsão.
A utilização de informação posterior ao jogo comprometeria a avaliação
de um sistema destinado à previsão pré-jogo.

A interpretação das previsões acrescenta outra dimensão ao problema.
Ouyang et al.~\cite{ouyang2024} combinam XGBoost e SHAP na análise de
resultados da NBA, utilizando dados das temporadas de 2021 a 2023.
Enquanto o primeiro produz previsões, o segundo estima a contribuição
das variáveis para essas previsões. A disponibilização de código pelos
autores constitui também um recurso para estudar a implementação.
Para um treinador, identificar os fatores associados a uma probabilidade
de vitória pode facilitar a interpretação do modelo. Essa explicação,
porém, não demonstra o efeito causal de alterar um alinhamento ou a
distribuição de minutos.

A previsão durante a partida utiliza informação diferente da previsão
pré-jogo. Kandhway~\cite{kandhway2025} analisa oito temporadas da NBA,
de 2015/2016 a 2022/2023, recorrendo a aprendizagem supervisionada e SHAP.
O estudo reporta aproximadamente 74\% de exatidão global, com valores
entre cerca de 63\% e 90\% em diferentes estados do jogo. A importância
das variáveis também se altera durante a partida. Estes resultados
fundamentam uma possível extensão dinâmica do projeto, mas não devem
ser comparados diretamente com previsões anteriores ao início do jogo,
que dispõem de menos informação.

A revisão sistemática de Li et al.~\cite{li2025}, que reúne 34 estudos
numa janela de pesquisa centrada em 2019--2024, ajuda a enquadrar estas
diferenças. Os autores salientam a preparação dos dados, a seleção de
características e a necessidade de melhorar a validação e diversificar
as competições estudadas. Valores elevados de exatidão, isoladamente,
não permitem ordenar os métodos quando variam os dados, as entradas e
os protocolos de avaliação. Para o presente trabalho, importa comparar
modelos nas mesmas condições temporais e analisar também a qualidade
das probabilidades que foram fornecidas ao planeador.
